% Môn: Vật lý
% Lớp: 11
% Chương: Chương IV. Dòng điện. Mạch điện
% Bài: L11C4  Bài 23. Điện trở. Định luật Ohm
% Số câu: 16
% App đọc file này trực tiếp — sửa rồi Commit trên GitHub, bấm Đồng bộ GitHub .tex

% L11C4  Bài 23. Điện trở. Định luật Ohm

% ===== Câu 1 =====
% ID: L11C4B23-01-TL
% Mức: TH
\dangbt{Chưa có dạng}
\begin{ex}
% ID: L11C4B23-01-TL
% Mức: TH
Hai điện trở R_1 và R_2 mắc nối tiếp thì điện trở tương đương lớn gấp 6,25 lần điện trở tương $R$ đương khi mắc song song. Tính tỉ số hai điện trở R_1. 2 $\nguon{ly11 kntt.pdf}$
\choiceTF

\loigiai{
Theo đầu bài ta có: $R_{1} + R_{2} = 6,25 \cdot \dfrac{R_{1}R_{2}}{R_{1} + R_{2}}$.

Đặt $x = \dfrac{R_{1}}{R_{2}}$ qua một số biến đổi

$\begin{matrix}
{\left( {R_{1} + R_{2}} \right)^{2} = 6,25 \cdot R_{1}R_{2}} \\
\left. \Leftrightarrow R_{1}^{2} + 2R_{1}R_{2} + R_{2}^{2} = 6,25R_{1}R_{2} \right. \\
\left. \Leftrightarrow\dfrac{R_{1}}{R_{2}} + 2 + \dfrac{R_{2}}{R_{1}} = 6,25 \right. \\
\left. \Leftrightarrow x + \dfrac{1}{x} - 4,25 = 0 \right.
\end{matrix}$

Ta thu được phương trình x^{2}-4,25x+1 = 0 .

Giải phương trình ta thu được $x_{1} = 4;x_{2} = \dfrac{1}{4}$.
}
\end{ex}

% ===== Câu 2 =====
% ID: L11C4B23-01-TN
% Mức: NB
\begin{ex}
% ID: L11C4B23-01-TN
% Mức: NB
Đơn vị đo điện trở là $\nguon{ly11 kntt.pdf}$
\choice
{\True ôm ($\Omega$).}
{fara (F).}
{henry (H).}
{oát (W).}
\loigiai{
Chọn A.

Đơn vị đo điện trở là ôm $(\Omega )$
}
\end{ex}

% ===== Câu 3 =====
% ID: L11C4B23-02-TN
% Mức: NB
\begin{ex}
% ID: L11C4B23-02-TN
% Mức: NB
Phát biểu nào sau đây sai. $\nguon{ly11 kntt.pdf}$
\choice
{Điện trở có vạch màu là căn cứ để xác định trị số.}
{Đối với điện trở nhiệt có hệ số dương, khi nhiệt độ tăng thì điện trở tăng.}
{\True Đối với điện trở biến đổi theo điện áp, khi $U$ tăng thì điện trở tăng.}
{Đối với điện trở quang, khi ánh sáng thích hợp rọi vào thì điện trở giảm.}
\loigiai{
Chọn C.

$C$ – sai
}
\end{ex}

% ===== Câu 4 =====
% ID: L11C4B23-03-TN
% Mức: NB
\begin{ex}
% ID: L11C4B23-03-TN
% Mức: NB
Đặc điểm của điện trở nhiệt có hệ số nhiệt điện trở $\nguon{ly11 kntt.pdf}$
\choice
{\True dương khi nhiệt độ tăng thì điện trở tăng.}
{dương khi nhiệt độ tăng thì điện trở giảm.}
{âm khi nhiệt độ tăng thì điện trở tăng.}
{âm khi nhiệt độ tăng thì điện trở giảm về bằng 0.}
\loigiai{
Chọn A.

Đặc điểm của điện trở nhiệt có hệ số nhiệt điện trở dương khi nhiệt độ tăng thì điện trở tăng. Công thức điện trở $R = \rho\dfrac{l}{S}$ nên khi nhiệt độ tăng thì điện trở suất tăng do đó điện trở tăng.
}
\end{ex}

% ===== Câu 5 =====
% ID: L11C4B23-04-TN
% Mức: NB
\begin{ex}
% ID: L11C4B23-04-TN
% Mức: NB
Nếu chiều dài và đường kính của một dây dẫn bằng đồng có tiết diện tròn được tăng lên gấp đôi thì điện trở của dây dẫn sẽ $\nguon{ly11 kntt.pdf}$
\choice
{không thay đổi.}
{tăng lên hai lần.}
{tăng lên gấp bốn lần.}
{\True giảm đi hai lần.}
\loigiai{
Chọn D.

Điện trở $R = \rho\dfrac{l}{S} = \rho\dfrac{l}{\dfrac{\pi d^{2}}{4}}$ nên khi chiều dài và đường kính của một dây dẫn bằng đồng có tiết diện tròn được tăng lên gấp đôi thì điện trở của dây dẫn sẽ giảm đi 2 lần.
}
\end{ex}

% ===== Câu 6 =====
% ID: L11C4B23-05-TN
% Mức: NB
\begin{ex}
% ID: L11C4B23-05-TN
% Mức: NB
Chọn biến đổi đúng trong các biến đổi sau. $\nguon{ly11 kntt.pdf}$
\choice
{1 $\Omega = 0$,001k $\Omega = 0$,0001 $M$ \Omega.}
{10 $\Omega = 0$,1k $\Omega = 0$,00001 $M$ \Omega.}
{1k $\Omega = 1000\Omega = 0$,01 $M$ \Omega.}
{\True 1 $M\Omega = 1000k\Omega = 1000000\Omega.$}
\loigiai{
Chọn D.

$1\,\mathrm{M}\Omega= 1000k\Omega= 1000000$ \Omega
}
\end{ex}

% ===== Câu 7 =====
% ID: L11C4B23-06-TN
% Mức: NB
\begin{ex}
% ID: L11C4B23-06-TN
% Mức: NB
Biến trở là $\nguon{ly11 kntt.pdf}$
\choice
{điện trở có thể thay đổi trị số và dùng để điều chỉnh chiều dòng điện trong mạch.}
{điện trở' có thể thay đổi trị số và dùng để điều chỉnh cường độ và chiều dòng điện trong mạch.}
{\True điện trở có thể thay đổi trị số và dùng để điều chỉnh cường độ dòng điện trong mạch.}
{điện trở không thay đổi trị số và dùng để điều chỉnh cường độ dòng điện trong mạch.}
\loigiai{
Chọn C.

Biến trở là điện trở có thể thay đổi trị số (tức là giá trị điện trở thay đổi) và dùng để điều chỉnh cường độ dòng điện trong mạch.
}
\end{ex}

% ===== Câu 8 =====
% ID: L11C4B23-07-TN
% Mức: NB
\begin{ex}
% ID: L11C4B23-07-TN
% Mức: NB
Trước khi mắc biến trở vào mạch điện để điều chỉnh cường độ dòng điện thì cần điều chỉnh biến trở có giá trị nào dưới đây? $\nguon{ly11 kntt.pdf}$
\choice
{Có giá trị bằng 0.}
{Có giá trị nhỏ.}
{Có giá trị lớn.}
{\True Có giá trị lớn nhất.}
\loigiai{
Chọn D.

Trước khi mắc biến trở vào mạch điện để điều chỉnh cường độ dòng điện thì cần điều chỉnh biến trở có giá trị lớn nhất, mục đích để cho dòng điện trong mạch có giá trị nhỏ nhất, sau đó giảm dần giá trị biến trở để tăng cường độ dòng điện.
}
\end{ex}

% ===== Câu 9 =====
% ID: L11C4B23-08-TN
% Mức: NB
\begin{ex}
% ID: L11C4B23-08-TN
% Mức: NB
Chọn phát biểu đúng về định luật Ohm. $\nguon{ly11 kntt.pdf}$
\choice
{Cường độ dòng điện chạy qua dây dẫn tỉ lệ với hiệu điện thế giữa hai đầu dây dẫn và điện trở của dây.}
{Cường độ dòng điện chạy qua dây dẫn tỉ lệ thuận với hiệu điện thế giữa hai đầu dây dẫn và không tỉ lệ với điện trở của dây.}
{\True Cường độ dòng điện chạy qua dây dẫn tỉ lệ thuận với hiệu điện thế giữa hai đầu dây dẫn và tỉ lệ nghịch với điện trở của dây.}
{Cường độ dòng điện chạy qua dây dẫn tỉ lệ nghịch với hiệu điện thế giữa hai đầu dây dẫn và tỉ lệ thuận với điện trở của dây.}
\loigiai{
Chọn C.

Định luật Ohm: Cường độ dòng điện chạy qua dây dẫn tỉ lệ thuận với hiệu điện thế giữa hai đầu dây dẫn và tỉ lệ nghịch với điện trở của dây.
}
\end{ex}

% ===== Câu 10 =====
% ID: L11C4B23-09-TN
% Mức: NB
\begin{ex}
% ID: L11C4B23-09-TN
% Mức: NB
Khi hiệu điện thế giữa hai đầu dây dẫn tăng thì $\nguon{ly11 kntt.pdf}$
\choice
{cường độ dòng điện chạy qua dây dẫn không thay đổi.}
{cường độ dòng điện chạy qua dây dẫn giảm, tỉ lệ với hiệu điện thế.}
{cường độ dòng điện chạy qua dây dẫn có lúc tăng, có lúc giảm.}
{\True cường độ dòng điện chạy qua dây dẫn tăng, tỉ lệ với hiệu điện thế.}
\loigiai{
Chọn D.

Khi hiệu điện thế giữa hai đầu dây dẫn tăng thì cường độ dòng điện chạy qua dây dẫn tăng, tỉ lệ với hiệu điện thế.
}
\end{ex}

% ===== Câu 11 =====
% ID: L11C4B23-10-TN
% Mức: NB
\begin{ex}
% ID: L11C4B23-10-TN
% Mức: NB
Cường độ dòng điện chạy qua dây dẫn tỉ lệ thuận với hiệu điện thế giữa hai đầu dây dẫn. Nếu tăng hiệu điện thế lên 1,6 lần thì $\nguon{ly11 kntt.pdf}$
\choice
{cường độ dòng điện tăng 3,2 lần.}
{cường độ dòng điện giảm 3,2 lần.}
{cường độ dòng điện giảm 1,6 lần.}
{\True cường độ dòng điện tăng 1,6 lần.}
\loigiai{
Chọn D.

Nếu tăng hiệu điện thế lên 1{,}6 lần thì cường độ dòng điện tăng 1{,}6 lần.
}
\end{ex}

% ===== Câu 12 =====
% ID: L11C4B23-11-TN
% Mức: NB
\begin{ex}
% ID: L11C4B23-11-TN
% Mức: NB
Cường độ dòng điện chạy qua điện trở $R = 5\Omega là 0$, $5\,\mathrm{A}$.Hiệu điện thế giữa hai đầu điện trở là $\nguon{ly11 kntt.pdf}$
\choice
{\True $2,5\,\mathrm{V}$.}
{$25\,\mathrm{V}$.}
{$0,5\,\mathrm{V}$.}
{$10\,\mathrm{V}$.}
\loigiai{
Chọn A.

Hiệu điện thế cần tìm $U= I.R= 0,5.5 = 2,5\,\mathrm{V}$
}
\end{ex}

% ===== Câu 13 =====
% ID: L11C4B23-12-TN
% Mức: NB
\begin{ex}
% ID: L11C4B23-12-TN
% Mức: NB
Đặt một hiệu điện thế $U = 12\$, $\mathrm{V}$ vào hai đầu một điện trở, cường độ dòng điện chạy qua điện trở là $2\,\mathrm{A}$.Nếu tăng hiệu điện thế lên hai lần thì cường độ dòng điện có giá trị $\nguon{ly11 kntt.pdf}$
\choice
{\True $4\,\mathrm{A}$.}
{$2\,\mathrm{A}$.}
{$1,2\,\mathrm{A}$.}
{$0,24\,\mathrm{A}$.}
\loigiai{
Chọn A.

Hiệu điện thế tăng 2 lần thì cường độ dòng điện tăng 2 lần.
}
\end{ex}

% ===== Câu 14 =====
% ID: L11C4B23-13-TN
% Mức: NB
\begin{ex}
% ID: L11C4B23-13-TN
% Mức: NB
Đặt vào hai đầu một điện trở $R$ một hiệu điện thế $U = 12\$, $\mathrm{V},$ cường độ dòng điện chạy qua điện trở là $1,5\,\mathrm{A}$. Nếu giữ nguyên hiệu điện thế nhưng muốn cường độ dòng điện chạy qua điện trở là giảm đi $0,5\,\mathrm{A}$ thì ta phải tăng giá trị điện trở thêm một lượng là $\nguon{ly11 kntt.pdf}$
\choice
{5,0 $\Omega$.}
{4,5 $\Omega$.}
{\True 4,0 $\Omega$.}
{5,5 $\Omega$.}
\loigiai{
Chọn C.

Điện trở ban đầu: $R_{1} = \dfrac{U_{1}}{I_{1}} = \dfrac{12}{1,5} = 8\Omega$

Điện trở lức sau: $R_{2} = \dfrac{U_{2}}{I_{2}} = \dfrac{U_{1}}{I_{1} - 0,5} = \dfrac{12}{1,5 - 0,5} = 12\Omega$

Điện trở tăng thêm 1 lượng 4 $\Omega$
}
\end{ex}

% ===== Câu 15 =====
% ID: L11C4B23-14-TN
% Mức: TH
\begin{ex}
% ID: L11C4B23-14-TN
% Mức: TH
Biểu thức đúng của định luật Ohm là $RU$ 1 $R\nguon{ly11 kntt.pdf}$
\choice
{$I = U$.}
{\True $I = R$.}
{$U = R$.}
{$U = I$.}
\loigiai{
Chọn B.

Biểu thức định luật Ohm: $I = \dfrac{U}{R}$
}
\end{ex}

% ===== Câu 16 =====
% ID: L11C4B23-15-TN
% Mức: TH
\begin{ex}
% ID: L11C4B23-15-TN
% Mức: TH
Muốn đo hiệu điện thế giữa hai cực của một nguồn điện, nhưng không có vôn kế, một học $\sin$ h đã sử dụng một ampe kế và một điện trở có giá trị $R = 50\Omega$ mắc nối tiếp nhau sau, đó mắc vào nguồn điện, biết ampe kế chỉ $1,2\,\mathrm{A}$.Hiệu điện thế giữa hai cực nguồn điện có giá trị bằng bao nhiêu? $\nguon{ly11 kntt.pdf}$
\choice
{$120\,\mathrm{V}$.}
{$50\,\mathrm{V}$.}
{$12\,\mathrm{V}$.}
{\True $60\,\mathrm{V}$.}
\loigiai{
Chọn D.

Hiệu điện thế $U= I.R=60\,\mathrm{V}$
}
\end{ex}
